\begin{theorem}[亚临界 mixed tensor 的真空律]\label{thm:conclusion-window6-subcritical-mixed-tensor-vacuum}
设
\[
V:=\RR^{X_6},
\qquad
V_{\CC}:=V\otimes_{\RR}\CC,
\qquad
\mathfrak g_6=\Lie\langle L_1,\dots,L_6\rangle=\mathfrak{sl}(V)\cong \mathfrak{sl}_{21}.
\]
则对任意非负整数 \(p,q\) 满足
\[
p+q<21,
\qquad
p\neq q,
\]
有
\[
\operatorname{Hom}_{\mathfrak g_6}\!\bigl(V^{\otimes p},V^{\otimes q}\bigr)=0.
\]
等价地，在总张量次数低于 \(21\) 的范围内，一切非平衡型 mixed tensor intertwiner 都严格消失。
\end{theorem}
\begin{proof}
若实空间
\[
\operatorname{Hom}_{\mathfrak g_6}\!\bigl(V^{\otimes p},V^{\otimes q}\bigr)
\]
非零，则其复化也非零；故只需证明
\[
\operatorname{Hom}_{\mathfrak{sl}_{21}}\!\bigl(V_{\CC}^{\otimes p},V_{\CC}^{\otimes q}\bigr)=0.
\]
由标准同构，
\[
\operatorname{Hom}_{\mathfrak{sl}_{21}}\!\bigl(V_{\CC}^{\otimes p},V_{\CC}^{\otimes q}\bigr)
\cong
\bigl((V_{\CC}^*)^{\otimes p}\otimes V_{\CC}^{\otimes q}\bigr)^{SL_{21}}.
\]
令 \(\zeta I\) 为 \(SL_{21}(\CC)\) 的中心元，其中 \(\zeta^{21}=1\)。它在
\[
(V_{\CC}^*)^{\otimes p}\otimes V_{\CC}^{\otimes q}
\]
上按标量 \(\zeta^{q-p}\) 作用。若上述不变量空间非零，则中心必须平凡作用，故对一切 \(\zeta^{21}=1\) 都有 \(\zeta^{q-p}=1\)，即
\[
21\mid (q-p).
\]
但由 \(p+q<21\) 可知 \(|q-p|<21\)，再结合 \(p\neq q\)，便不可能有 \(21\mid(q-p)\)。故不变量空间为零，从而原实空间亦为零。证毕。
\end{proof}

\begin{theorem}[亚临界 Schur--Weyl 刚性]\label{thm:conclusion-window6-subcritical-schur-weyl-rigidity}
对每个整数
\[
0\le p<21,
\]
有精确同构
\[
\operatorname{End}_{\mathfrak g_6}(V^{\otimes p})
\cong
\RR[\mathfrak S_p]_{\mathrm{place}},
\]
即一切 \(\mathfrak g_6\)-等变自同态都恰由 tensor slots 的置换线性张成。换言之，在稳定范围 \(p<21\) 内，window-\(6\) 的完整 edge--flux 动力学在张量层面不容纳任何超出置换对称的额外中心化子。
\end{theorem}
\begin{proof}
复化后只需证明
\[
\operatorname{End}_{\mathfrak{sl}_{21}}(V_{\CC}^{\otimes p})
\cong
\CC[\mathfrak S_p]_{\mathrm{place}}.
\]
由定理 \ref{thm:window6-lie-envelope-sl21}，\(V_{\CC}\) 是 \(\mathfrak{sl}_{21}\) 的标准模。于是经典 Schur--Weyl 对偶给出
\[
\operatorname{End}_{SL_{21}}(V_{\CC}^{\otimes p})
\cong
\CC[\mathfrak S_p]_{\mathrm{place}}.
\]
在此多项式表示范畴中，\(\mathfrak{sl}_{21}\)-等变与 \(SL_{21}\)-等变一致，因此得到复系数结论。由于 place permutations 全由实矩阵给出，取实型即得
\[
\operatorname{End}_{\mathfrak g_6}(V^{\otimes p})
\cong
\RR[\mathfrak S_p]_{\mathrm{place}}.
\]
证毕。
\end{proof}

\begin{theorem}[临界次数之前的不变量真空与 \texorpdfstring{$21$}{21} 阶 Pl\"ucker 生成]\label{thm:conclusion-window6-degree21-plucker-first-invariants}
对任意整数 \(M\ge 1\)，令 \(\mathfrak g_6\) 在 \(V^{\oplus M}\) 上按对角方式作用。则
\[
\RR[V^{\oplus M}]^{\mathfrak g_6}_d=0
\qquad
(1\le d<21),
\]
而在临界阶 \(d=21\) 恰有
\[
\RR[V^{\oplus M}]^{\mathfrak g_6}_{21}
\cong
\Lambda^{21}(\RR^M).
\]
更具体地，当 \(M\ge 21\) 时，全部 \(21\) 次不变量由 \(21\times 21\) 子式线性张成；当 \(M<21\) 时，甚至 \(21\) 次不变量也完全消失。
\end{theorem}
\begin{proof}
复化后，只需证明
\[
\CC[V_{\CC}^{\oplus M}]^{\mathfrak{sl}_{21}}_d=0
\qquad
(1\le d<21),
\]
且
\[
\CC[V_{\CC}^{\oplus M}]^{\mathfrak{sl}_{21}}_{21}
\cong
\Lambda^{21}(\CC^M).
\]
将 \(V_{\CC}^{\oplus M}\) 视为 \(21\times M\) 矩阵的列空间，经典 \(SL_n\) 第一基本定理断言：\(SL_{21}\) 在 \(V_{\CC}^{\oplus M}\) 上的多项式不变量环由全部 \(21\times 21\) 子式生成。故在次数 \(d<21\) 时，不存在非零齐次不变量；而在次数 \(21\) 时，不变量空间正由这些子式线性张成，其系数空间自然同构于
\[
\Lambda^{21}(\CC^M).
\]
取实型即得所述结论。证毕。
\end{proof}

\begin{theorem}[rank-one routing 的完备可综合性]\label{thm:conclusion-window6-rankone-routing-completeness}
\leanverified{paper\_conclusion\_window6\_rankone\_routing\_completeness}
对任意有序对
\[
(u,v)\in X_6\times X_6,
\]
存在带常数项的六变量非交换实多项式
\[
p_{vu}\in \RR\langle x_1,\dots,x_6\rangle
\]
使得
\[
p_{vu}(L_1,\dots,L_6)=E_{vu},
\]
其中 \(E_{vu}\) 为标准矩阵单位：
\[
E_{vu}(e_u)=e_v,
\qquad
E_{vu}(e_w)=0\quad (w\neq u).
\]
因此，对任意由坐标子集张成的子空间
\[
U=\operatorname{span}\{e_u:u\in Y\},
\qquad
W=\operatorname{span}\{e_v:v\in Z\},
\]
以及任意线性映射 \(T:U\to W\)，都存在非交换多项式 \(p_T\) 使得
\[
p_T(L_1,\dots,L_6)|_U=T,
\qquad
\operatorname{im}p_T(L_1,\dots,L_6)\subseteq W.
\]
\end{theorem}
\begin{proof}
由定理 \ref{thm:window6-edge-flux-full-matrix-saturation}，
\[
\RR\langle I,L_1,\dots,L_6\rangle=\operatorname{End}(V).
\]
故每个矩阵单位 \(E_{vu}\in \operatorname{End}(V)\) 都必可表示为某个带常数项的非交换多项式 \(p_{vu}(L_1,\dots,L_6)\)。这证明了第一部分。

对第二部分，只需把 \(T\) 延拓为某个 \(\widetilde T\in \operatorname{End}(V)\)，使其在 \(U\) 上等于 \(T\)，在 \(U\) 的坐标补空间上取零，且像包含于 \(W\)。于是
\[
\widetilde T=\sum_{u\in Y,\ v\in Z}t_{vu}E_{vu}.
\]
逐项代入第一部分的表示式并线性组合，即得某个 \(p_T\) 满足
\[
p_T(L_1,\dots,L_6)=\widetilde T.
\]
从而其在 \(U\) 上的限制恰为 \(T\)。证毕。
\end{proof}

\begin{theorem}[投影阴影布尔代数的精确综合与中心坍缩]\label{thm:conclusion-window6-boolean-shadow-center-collapse}
定义
\[
\mathscr B_6:=\{P_Y:Y\subseteq X_6\}\subset \operatorname{End}(V),
\]
其中 \(P_Y\) 为坐标投影。则：
\begin{enumerate}
  \item \(\mathscr B_6\subset \RR\langle I,L_1,\dots,L_6\rangle\)，亦即每个 \(P_Y\) 都可由 \(L_1,\dots,L_6\) 的非交换多项式精确综合；
  \item 映射
  \[
  Y\longmapsto P_Y
  \]
  把 \(\mathcal P(X_6)\) 同构到 \(\operatorname{End}(V)\) 内的一个布尔代数，并满足
  \[
  P_{Y\cap Z}=P_YP_Z,
  \qquad
  P_{Y\cup Z}=P_Y+P_Z-P_YP_Z,
  \qquad
  P_{Y^c}=I-P_Y;
  \]
  \item 该布尔代数的动力学中心严格坍缩：
  \[
  \mathscr B_6\cap \operatorname{Comm}(L_1,\dots,L_6)=\{0,I\}.
  \]
\end{enumerate}
因此，window-\(6\) 同时具有全布尔精确读写能力与零非平凡中央阴影。
\end{theorem}
\begin{proof}
第一条由推论 \ref{cor:window6-exact-projector-synthesis} 直接得到。第二条只是坐标投影的标准恒等式，因此
\[
Y\mapsto P_Y
\]
把幂集布尔代数同构到 diagonal idempotents 的布尔代数。

对第三条，若 \(P_Y\in \mathscr B_6\) 与全部 \(L_i\) 对易，则由定理 \ref{thm:window6-edge-flux-commutant-rigidity}，
\[
P_Y\in \operatorname{Comm}(L_1,\dots,L_6)=\RR I.
\]
但 \(P_Y\) 还是幂等元，故作为标量只能取 \(0\) 或 \(1\)。于是
\[
P_Y\in\{0,I\},
\]
等价于 \(Y=\varnothing\) 或 \(Y=X_6\)。证毕。
\end{proof}

\begin{corollary}[可编程性与守恒性的严格分离]\label{cor:conclusion-window6-programmable-not-conserved}
设 \(\mathfrak G_6\) 为由 boundary/cyclic 分层、shell 分层、\(L/P^\circ/P/\Sigma\) 型纤维子簇以及它们的有限布尔组合所生成的几何扇区族。则对每个 \(\mathcal Y\in \mathfrak G_6\)，其对应坐标子空间
\[
V_{\mathcal Y}:=\operatorname{span}\{e_w:w\in \mathcal Y\}
\]
满足如下严格二分：
\begin{enumerate}
  \item \(V_{\mathcal Y}\) 可由 edge--flux 非交换多项式精确读写与精确路由；
  \item 除 \(0\) 与 \(V\) 外，\(V_{\mathcal Y}\) 从不构成公共不变子空间，更不构成守恒扇区。
\end{enumerate}
并且对任意两个非零几何扇区 \(\mathcal Y,\mathcal Z\in \mathfrak G_6\) 以及任意
\[
1\le r\le \min\{\dim V_{\mathcal Y},\dim V_{\mathcal Z}\},
\]
都存在一个 edge--flux 多项式 \(p_{\mathcal Z\leftarrow \mathcal Y}^{(r)}\) 实现秩 \(r\) 转导
\[
V_{\mathcal Y}\to V_{\mathcal Z}.
\]
因此，window-\(6\) 中
\[
\boxed{\text{exactly programmable}}
\qquad\not\Rightarrow\qquad
\boxed{\text{dynamically conserved}}.
\]
\end{corollary}
\begin{proof}
由定理 \ref{thm:conclusion-window6-boolean-shadow-center-collapse}，全部几何扇区对应的坐标投影都可精确综合，故精确读写成立。再由定理 \ref{thm:conclusion-window6-rankone-routing-completeness}，任意两个这类坐标子空间之间的任意秩 \(r\) 线性映射都可由 edge--flux 多项式实现，故精确路由成立。

另一方面，若 \(V_{\mathcal Y}\) 是全部 \(L_i\) 的公共不变子空间，则其投影 \(P_{\mathcal Y}\) 必与全部 \(L_i\) 对易；由定理 \ref{thm:conclusion-window6-boolean-shadow-center-collapse}，只可能有
\[
P_{\mathcal Y}\in\{0,I\},
\]
即 \(\mathcal Y=\varnothing\) 或 \(\mathcal Y=X_6\)。故非平凡几何扇区永远不是守恒扇区。证毕。
\end{proof}

\begin{theorem}[window-\texorpdfstring{$6$}{6} 的账本维数与 Lie 闭包维数同一]\label{thm:conclusion-window6-ledger-lie-dimension-identity}
\leanverified{paper\_conclusion\_window6\_ledger\_lie\_dimension\_identity}
记
\[
V:=\RR^{X_6},
\qquad
\mathfrak k_6:=\Lie\langle [L_i,L_j]:1\le i<j\le 6\rangle,
\qquad
\mathfrak g_6:=\Lie\langle L_1,\dots,L_6\rangle.
\]
则有严格同一
\[
|X_6|
=
\dim_{\FF_2}\bigl((G_6^{\mathrm{bin}})^{\mathrm{ab}}\bigr)
=
\dim_{\RR}V
=21,
\]
并且
\[
\mathfrak k_6\cong \mathfrak{so}(21),
\qquad
\mathfrak g_6\cong \mathfrak{sl}(21).
\]
因此，window-\(6\) 的最小完备异常账本维数与最大正交闭包的标准忠实实表示维数完全一致，并共同锁定为 \(21\)。
\end{theorem}
\begin{proof}
由定义
\[
V=\RR^{X_6}
\]
立得
\[
\dim_{\RR}V=|X_6|.
\]
而 window-\(6\) 满足
\[
|X_6|=21.
\]
另一方面，定理 \ref{thm:conclusion-window6-static-anomaly-ledger-dynamic-budget-bifurcation} 已给出
\[
\dim_{\FF_2}\bigl((G_6^{\mathrm{bin}})^{\mathrm{ab}}\bigr)=21.
\]
故三者同一
\[
|X_6|
=
\dim_{\FF_2}\bigl((G_6^{\mathrm{bin}})^{\mathrm{ab}}\bigr)
=
\dim_{\RR}V
=21
\]
成立。

再由定理 \ref{thm:conclusion-window6-observable-sectors-vs-cartan-saturated-ambient}，
\[
\mathfrak k_6=\mathfrak{so}(V),
\qquad
\mathfrak g_6=\mathfrak{sl}(V).
\]
结合 \(\dim_{\RR}V=21\)，遂得
\[
\mathfrak k_6\cong \mathfrak{so}(21),
\qquad
\mathfrak g_6\cong \mathfrak{sl}(21).
\]
最后，\(\mathfrak{so}(21)\) 的标准 \(21\) 维向量表示显然忠实，而型 \(B_{10}\) 的经典表示论表明不存在更低维的非平凡忠实实表示；故 window-\(6\) 中最小完备异常账本维数与最大正交闭包的标准忠实实表示维数完全重合。证毕。
\end{proof}

\begin{conjecture}[满 edge--flux 饱和窗口的临界维数原理]\label{conj:conclusion-windowm-critical-dimension-principle}
设一般分辨率 \(m\) 满足
\[
\mathcal B_m:=\RR\langle I,L_1^{(m)},\dots,L_r^{(m)}\rangle=\operatorname{End}(V_m),
\qquad
\dim V_m=N_m,
\]
\[
\operatorname{Comm}(L_1^{(m)},\dots,L_r^{(m)})=\RR I,
\qquad
V_{m,\CC}\ \text{是 }\mathfrak{sl}_{N_m}\text{ 的标准模}.
\]
则应当存在一条仅由临界维数 \(N_m\) 支配的统一阈值律：
\begin{enumerate}
  \item 若 \(p+q<N_m\) 且 \(p\neq q\)，则
  \[
  \operatorname{Hom}_{\mathfrak g_m}\!\bigl(V_m^{\otimes p},V_m^{\otimes q}\bigr)=0;
  \]
  \item 若 \(p<N_m\)，则
  \[
  \operatorname{End}_{\mathfrak g_m}(V_m^{\otimes p})
  \cong
  \RR[\mathfrak S_p]_{\mathrm{place}};
  \]
  \item 对任意 \(M\ge 1\)，
  \[
  \RR[V_m^{\oplus M}]_d^{\mathfrak g_m}=0
  \qquad
  (1\le d<N_m),
  \]
  而次数 \(N_m\) 的不变量首次由 \(N_m\times N_m\) 子式生成；
  \item 全部 projector-shadows 组成一个可精确综合的布尔代数，但其动力学中心仅含平凡投影。
\end{enumerate}
window-\(6\) 正是这一原理的首个完整模型，其中
\[
N_6=|X_6|=21.
\]
\end{conjecture}
